%==================================================================================================================
%======================================Figure photovoltaïque=======================================================
%==================================================================================================================
\newcommand{\figschemaphotovoltaique}{%

\begin{tikzpicture}[x=1cm,y=1cm,scale=0.84,transform shape,>=Stealth]

    % Soleil
    \fill[pvSun!38] (0.6,4.0) circle (0.48);
    \fill[pvSun!90] (0.6,4.0) circle (0.33);

    % Couches : proportions des blocs
    \pvblock{2.10}{1.20}{0.58}{2.00}{}{pvITO}
    \pvblock{2.74}{1.28}{1.18}{2.00}{}{pvCIGS}
    \pvblock{4.02}{1.36}{0.34}{2.00}{}{pvZnO}
        \pvblock{4.46}{1.44}{0.38}{2.00}{}{pvCdS}
    \pvblock{4.94}{1.52}{0.58}{2.00}{}{pvMetal}

    % Labels des blocs
    \node[font=\scriptsize\bfseries,text=yellow!90!orange] at (2.39,2.20) {ITO};
    \node[font=\scriptsize\bfseries,text=yellow!90!orange] at (3.33,2.28) {CIGS};
    \node[font=\tiny\bfseries,text=yellow!90!orange,rotate=90] at (4.19,2.36) {ZnO};
    \node[font=\tiny\bfseries,text=yellow!90!orange,rotate=90] at (4.65,2.44) {CdS};
    \node[font=\scriptsize\bfseries,text=yellow!90!orange] at (5.23,2.52) {Métal};

    % Flux solaire
    \draw[->,draw=pvSun,line width=1.2pt]
      (1.10,3.25) .. controls (1.35,2.80) and (1.65,2.45) .. (2.00,2.20);

    \node[align=center,font=\scriptsize\bfseries,text=pvFlux]
      at (0.8,2.8) {Flux\\solaire};

    % e- depuis l'anode
    \draw[->,draw=pvSun,line width=1.2pt]
      (0.95,1.10) .. controls (1.30,0.95) and (1.60,1.00) .. (2.00,1.18);

    \node[align=center,font=\scriptsize\bfseries,text=pvFlux]
      at (0.82,1.55) {$e^-$ depuis\\l'anode};

    % e- vers la catalyse
    \draw[->,draw=pvSun,line width=1.2pt]
      (5.90,2.50) .. controls (6,2.5) and (6.3,2.3) .. (7,2.95);

    \node[align=center,font=\scriptsize\bfseries,text=pvFlux]
      at (6.35,2.1) {$e^-$ vers\\catalyse};

  \end{tikzpicture}
}
%==================================================================================================================
%======================================Fin Figure photovoltaïque===================================================
%==================================================================================================================

%------------------------------------------------------------------------------------------------------------------

%==================================================================================================================
%======================================Figure électrochimique=======================================================
%==================================================================================================================
\newcommand{\figschemaelectrochimique}{%

\begin{tikzpicture}[x=1cm,y=1cm,scale=0.7,transform shape,>=Stealth,
                    font=\scriptsize\bfseries]
 
  % Métal (électrode gauche)
  \fill[pvMetal] (0,0) rectangle (0.7,3);
  \node[text=white,rotate=90] at (0.35,1.5) {Métal};
 
  % Couche poreuse + catalyseur
  \fill[pvCdS] (0.7,0) rectangle (2.0,3);
  \foreach \y in {0.4,1.0,1.6,2.2,2.7}{
    \draw[pvSide!50,line width=0.4pt]
      (0.75,\y) .. controls (1.2,\y+0.25) and (1.5,\y-0.25) .. (1.95,\y);
  }
  \foreach \p in {(1.0,0.7),(1.5,1.3),(1.1,1.9),(1.6,2.4),(1.3,0.4)}{
    \fill[pvSide] \p circle (0.13);
    \node[text=white,font=\tiny] at \p {Cat};
  }
 
  % Électrolyte (cuve centrale)
  \fill[pvITO] (2.0,0) rectangle (4.6,3);
  \draw[pcNavy!40,line width=0.6pt] (2.0,0) rectangle (4.6,3);
  \node[text=pcNavy] at (3.3,1.7) {Électrolyte};
 
  % Membrane (droite)
  \fill[pvMetal] (4.6,0) rectangle (5.3,3);
  \node[text=white,rotate=90] at (4.95,1.5) {Membrane};
 
  % Flèche électrons : label au-dessus, flèche courbe vers le métal
  \node[align=center,text=pcNavy,font=\scriptsize\bfseries] at (-0.4,3.5)
    {$e^-$ depuis la partie PV};
  \draw[->,draw=pcNavy,line width=1.4pt]
    (-1.0,3.05) .. controls (-0.8,2.5) and (-0.5,2.3) .. (0.05,2.2);
 
  % Flèche protons : label au-dessus, flèche courbe vers la membrane
  \node[align=center,text=pcNavy,font=\scriptsize\bfseries] at (5.7,3.5)
    {Flux de protons};
  \draw[->,draw=pcNavy,line width=1.4pt]
    (6.3,3.05) .. controls (6.1,2.5) and (5.8,2.3) .. (5.25,2.2);
 
\end{tikzpicture}

}
%==================================================================================================================
%======================================Fin Figure électrochimique==================================================
%==================================================================================================================

%------------------------------------------------------------------------------------------------------------------

%==================================================================================================================
%======================================Figure démarche modèle======================================================
%==================================================================================================================
\newcommand{\figschemademarchemodelo}{%
    \begin{tikzpicture}[
      >=Stealth,
      node distance=0.8cm,
      objective/.style={
        draw=pcSky,
        fill=pcSkyLight,
        rounded corners=7pt,
        minimum width=2.6cm,
        minimum height=1.1cm,
        align=center,
        text=pcNavy,
        font=\footnotesize\bfseries
      },
      arrow/.style={
        ->,
        thick,
        draw=pcBlue
      }
    ]

      \node[objective] (mod) {Modéliser\\les sous-parties};

      \node[objective, right=of mod] (val)
        {Valider\\avec l'expérimental};

      \node[objective, right=of val] (coup)
        {Coupler\\les modèles};

      \node[objective, right=of coup] (opt)
        {Optimiser\\la photocathode};

      \draw[arrow] (mod) -- (val);
      \draw[arrow] (val) -- (coup);
      \draw[arrow] (coup) -- (opt);

    \end{tikzpicture}
}
%==================================================================================================================
%======================================Fin Figure démarche modèle==================================================
%==================================================================================================================


%==================================================================================================================
%======================================Recombinaisonss==================================================
%==================================================================================================================


\newcommand{\figrecombinaison}{%
\begin{tikzpicture}[
    >=Latex, font=\small,
    band/.style={draw=pcInk!45},
    elec/.style={circle, fill=pcBlue, draw=pcNavy, line width=0.4pt, inner sep=1.7pt},
    hole/.style={circle, draw=pcSun, fill=white, line width=0.8pt, inner sep=1.4pt},
    trans/.style={->, thick, pcNavy},
]
\def\W{3.0}\def\yc{2.4}\def\yv{0.6}\def\bh{0.55}
\def\recbandes{%
  \fill[pcSkyLight] (0,\yc) rectangle (\W,\yc+\bh);
  \fill[pcSun!15]    (0,\yv-\bh) rectangle (\W,\yv);
  \draw[band] (0,\yc) rectangle (\W,\yc+\bh);
  \draw[band] (0,\yv-\bh) rectangle (\W,\yv);
  \node[font=\scriptsize\bfseries, pcNavy]        at (\W/2,\yc+\bh/2) {BC};
  \node[font=\scriptsize\bfseries, pcSun!75!black] at (\W/2,\yv-\bh/2) {BV};
}
% 1. RADIATIVE
\begin{scope}[xshift=0cm]
  \recbandes
  \node[left, font=\scriptsize, pcInk] at (0,\yc) {$E_c$};
  \node[left, font=\scriptsize, pcInk] at (0,\yv) {$E_v$};
  \node[elec] at (0.8,\yc) {}; \node[hole] at (0.8,\yv) {};
  \draw[trans] (0.8,\yc-0.03) -- (0.8,\yv+0.03);
  \draw[->, pcSunDeep, thick, decorate,
        decoration={snake, amplitude=1pt, segment length=6pt, post length=3pt}]
        (1.05,1.5) -- (2.5,1.5);
  \node[pcSunDeep, font=\footnotesize] at (2.1,1.85) {$h\nu$};
  \node[font=\footnotesize\bfseries, pcInk] at (\W/2,-0.85) {Radiative};
  \node[font=\scriptsize, align=center, text=pcInk!70] at (\W/2,-1.35) {émission\\ d'un photon};
\end{scope}
% 2. SRH
\begin{scope}[xshift=4.3cm]
  \recbandes
  \draw[thick, pcInk] (0.85,1.5) -- (2.15,1.5);
  \node[font=\scriptsize, pcInk] at (\W/2,1.75) {$E_t$};
  \node[elec] at (1.0,\yc) {}; \node[hole] at (1.0,\yv) {};
  \draw[trans] (1.0,\yc-0.03) -- (1.05,1.55);
  \draw[trans] (1.15,1.45) -- (1.0,\yv+0.03);
  \node[font=\scriptsize, text=pcInk!70] at (2.35,1.0) {phonons};
  \node[font=\footnotesize\bfseries, pcInk] at (\W/2,-0.85) {SRH};
  \node[font=\scriptsize, align=center, text=pcInk!70] at (\W/2,-1.35) {assistée par\\ un défaut};
\end{scope}
% 3. AUGER
\begin{scope}[xshift=8.6cm]
  \recbandes
  \node[elec] at (0.7,\yc) {}; \node[hole] at (0.7,\yv) {};
  \draw[trans] (0.7,\yc-0.03) -- (0.7,\yv+0.03);
  \node[elec] at (2.15,\yc) {};
  \draw[trans] (2.15,\yc+0.03) -- (2.15,\yc+0.95);
  \node[font=\scriptsize, align=center, text=pcInk!70] at (2.15,\yc+1.25) {3\textsuperscript{e} porteur};
  \draw[->, dashed, pcInk!55] (0.9,1.5) .. controls (1.5,2.05) .. (2.1,\yc+0.35);
  \node[font=\footnotesize\bfseries, pcInk] at (\W/2,-0.85) {Auger};
  \node[font=\scriptsize, align=center, text=pcInk!70] at (\W/2,-1.35) {énergie à un\\ 3\textsuperscript{e} porteur};
\end{scope}
% Légende
\node[elec] at (2.0,-2.15) {};
\node[right, font=\scriptsize, pcInk] at (2.15,-2.15) {électron ($e^-$)};
\node[hole] at (5.4,-2.15) {};
\node[right, font=\scriptsize, pcInk] at (5.55,-2.15) {trou ($h^+$)};
\end{tikzpicture}%
}